\section{Point-in-Time Financial and Calibrated Mechanism Records}
\label{sec:supp-point-in-time-financial}

This section supplies the extended records for two distinct evidence layers.
The \emph{Point-in-Time Portfolio-Library Retention Study} is retrospective
financial evidence from licensed data. The \emph{Calibrated Closure-Option
Mechanism Experiment} is design-locked synthetic evidence whose nuisance law
uses public-safe aggregate calibration. The first is the article's financial
case study; the second isolates the bridge mechanism. Neither is an alpha,
forecasting, causal, or deployable-performance experiment.

\subsection{Point-in-time portfolio-library retention study}

\subsubsection{Chronology and information firewall}

For a compression origin \(y\), formation data from \(y-5\) through \(y-3\)
construct the source library and its operating frontier. The compression block
from \(y-2\) through \(y\) estimates strategy profiles and solves the safe and
frontier-only arms. Proposal year \(y+1\) enumerates every grammar-defined
innovation option enabled by each retained library and freezes the protected
policy choices. Evaluation year \(y+2\) is opened only after those choices are
hash-sealed. Proposal values never enter compression; evaluation values never
enter compression, option enumeration, uncertainty estimation, or adoption.

The primary universe is point-in-time liquid common equity. Plain ETFs form a
separate registered replication and are never pooled with common equities.
Eligibility does not require later survival. Observed delisting returns are
retained, terminal positions move to cash, and future completeness is not an
eligibility screen. One ETF origin fails its predeclared minimum-eligibility
gate and remains unavailable in every denominator rather than being silently
removed.

The source grammar contains portfolio-level signals, filters, holding horizons,
cross-sectional allocation, gross exposure, and exit rules. Capability closure
contains executable rule and interface requirements plus point-in-time data,
corporate-action and delisting, portfolio-aggregation, and cost/capacity
artifacts. The public manifest records \AORPointInTimeCapabilities\ complete
capabilities in every gate-passed cell. These are declared model artifacts;
they do not prove that a real institution's undocumented code or tacit
knowledge can be recovered from labels alone.

\subsubsection{Compression evidence and exact rechecks}

The innovation-safe arm preserves the entire source frontier and complete
declared closure. The comparator preserves only the frontier and is augmented
under the safe arm's realized burden cap with a frozen preproposal order. Its
action set is nested inside the safe action set before any proposal score is
read. Cash is always feasible, and the safe policy inherits the exact
comparator choice unless a simultaneous moving-block lower bound for an added
option clears the registered hurdle.

\AORPointInTimePreprocessingSolutions\ gate-passed safe cells reduce to an
exact preprocessing certificate. The remaining
\AORPointInTimeSolverSolutions\ are returned by HiGHS with
\texttt{OPTIMAL} termination. Every lifted selection then passes independent
original-instance checks for mandatory inactive retention, exact frontier
equality, exact identity-closure equality, and exact integer burden. These
checks certify the returned library's feasibility and burden. They do not turn
the solver's global-optimality status into Lean or exhaustive proof.

% Generated by julia/scripts/build_aor_manuscript_evidence_inputs.jl.
% Committed machine sources: experiments/financial_strategy_library_panel_v3/predecision_computation/PREDECISION_COMPUTATION_MANIFEST.toml; experiments/financial_strategy_library_panel_v3/evaluation_results/EVALUATION_RESULT_MANIFEST.toml; experiments/financial_strategy_library_panel_v3/PREDECISION_COMPUTATION_RESULT_SEAL.toml; experiments/financial_strategy_library_panel_v3/EVALUATION_RESULT_SEAL.toml.
\begin{table}[tbp]
\centering\small
\setlength{\tabcolsep}{3.6pt}
\resizebox{\textwidth}{!}{%
\begin{tabular}{@{}lrrrrrr@{}}
\toprule
Universe & Origins & Mean source burden & Mean safe burden & Mean reduction & Added options & Held-out CE difference \\
\midrule
Common equity & 19 & 8169.2 & 2183.8 & 73.1\% & 65.95 & 0.0000 \\
Plain ETF & 18 & 8909.2 & 2231.7 & 74.7\% & 66.06 & -0.0022 \\
\bottomrule
\end{tabular}
}
\caption{Point-in-time portfolio-library retention results. Burden is the registered validation-computation index, not currency. Added options are safe-library proposal actions absent from the budget-matched frontier-only library. The held-out column is innovation-safe minus frontier-only certainty-equivalent difference: all 19 common-equity choices coincide, while one of 18 complete ETF choices differs. Safe solutions pass exact original-instance frontier, closure, inactive-retention, and burden rechecks; solver-reported optimality remains solver evidence.}
\label{tab:point-in-time-financial-study}
\end{table}


The burden reductions in the table are origin-level means of
\(1-W(L^{\rm safe})/W(L)\). Burden is a prespecified
validation-computation index, not observed money, staff time, or compute use.
The option count is the number of safe-library proposal actions unavailable to
the budget-matched frontier-only library. It measures retained opportunity,
not realized performance.

\subsubsection{Origin-level held-out outcomes}

The common-equity protected policy never exercises an added innovation option,
so all held-out safe-minus-frontier-only contrasts are mechanically zero. In
the ETF replication, one protected replacement is exercised and realizes a
negative held-out contrast. The failure is retained; no alternative selector
is promoted after evaluation.

% Generated by julia/scripts/build_aor_manuscript_evidence_inputs.jl.
% Committed machine sources: experiments/financial_strategy_library_panel_v3/evaluation_results/EVALUATION_RESULT_MANIFEST.toml.
\begin{longtable}{@{}lrrrr@{}}
\caption{Origin-level point-in-time held-out contrasts. CE differences are innovation-safe minus frontier-only. A dash marks the prespecified failed ETF universe gate; it is retained, not imputed.}\label{tab:point-in-time-origin-results}\\
\toprule
Origin & Equity CE diff. & Same choice & ETF CE diff. & Same choice \\
\midrule
\endfirsthead
\toprule
Origin & Equity CE diff. & Same choice & ETF CE diff. & Same choice \\
\midrule
\endhead
2005 & 0.0000 & yes & -- & -- \\
2006 & 0.0000 & yes & 0.0000 & yes \\
2007 & 0.0000 & yes & 0.0000 & yes \\
2008 & 0.0000 & yes & 0.0000 & yes \\
2009 & 0.0000 & yes & 0.0000 & yes \\
2010 & 0.0000 & yes & 0.0000 & yes \\
2011 & 0.0000 & yes & 0.0000 & yes \\
2012 & 0.0000 & yes & 0.0000 & yes \\
2013 & 0.0000 & yes & 0.0000 & yes \\
2014 & 0.0000 & yes & 0.0000 & yes \\
2015 & 0.0000 & yes & 0.0000 & yes \\
2016 & 0.0000 & yes & -0.0390 & no \\
2017 & 0.0000 & yes & 0.0000 & yes \\
2018 & 0.0000 & yes & 0.0000 & yes \\
2019 & 0.0000 & yes & 0.0000 & yes \\
2020 & 0.0000 & yes & 0.0000 & yes \\
2021 & 0.0000 & yes & 0.0000 & yes \\
2022 & 0.0000 & yes & 0.0000 & yes \\
2023 & 0.0000 & yes & 0.0000 & yes \\
\bottomrule
\end{longtable}


The origin is the financial analysis unit. Strategy specifications and daily
observations are not treated as independent replications. A zero contrast from
identical actions does not establish economic equivalence of the two libraries,
and a negative realized contrast at one origin does not identify the population
value of closure.

\subsubsection{Registered robustness and diagnostic boundaries}

The proposal-stage robustness grid varies transaction costs, risk aversion,
and adoption hurdles without reopening compression or using evaluation data.
Common-equity safe and comparator choices coincide throughout. ETF differences
remain confined to the same origin and vanish at the higher cost settings.
Leave-one-origin-out summaries perform aggregation only, with no reselection.
A larger universe cap and a permuted-closure sham that were not frozen before
proposal outcomes became known remain explicitly unavailable.

The registered forced-maximum policy omits the protected cash/adoption rule and
is retained only as a negative control. Full-family best-candidate, White
reality-check, SPA, Deflated-Sharpe, and PBO summaries are search diagnostics;
they never feed back into a policy choice. Capacity records likewise score
frozen policies without reselection over the declared spread, impact, AUM, and
participation-cap grid. Missing capacity inputs, incomplete return paths, and
hard participation-cap breaches remain failures rather than dropped rows.

\subsection{Calibrated closure-option mechanism experiment}

\subsubsection{Calibration and registered regimes}

The mechanism experiment begins from two equal-burden libraries with equal
current operating frontiers. Only the innovation-safe library retains the
bridge capability for a declared descendant project. Its net bridge margin is
\[
M=\beta^d\rho^d\pi G-\kappa,
\]
and the oracle adopts exactly when \(M>0\). The frontier-only project is
closure-infeasible.

The noise calibration uses every complete proposal candidate path in the
gate-passed point-in-time cells: \AORPointInTimeCompletePaths\ paths and
\AORPointInTimeResidualObservations\ residual observations. Cross-strategy
date means are removed within cells before estimating scale, lag-one
dependence, and tail thickness. No sign or treatment margin is selected from
financial outcomes. An independent null simulation fixes the learner's
one-sided threshold; disjoint seeds generate test worlds; and the powered
margin is fixed prospectively at twice that threshold.

% Generated by julia/scripts/build_aor_manuscript_evidence_inputs.jl.
% Committed machine sources: experiments/financial_strategy_library_panel_v4/results/RESULT_MANIFEST.toml; experiments/financial_strategy_library_panel_v4/RESULT_SEAL.toml.
\begin{table}[tbp]
\centering\small
\setlength{\tabcolsep}{3.5pt}
\resizebox{\textwidth}{!}{%
\begin{tabular}{@{}lrrrr@{}}
\toprule
Regime & True margin & Oracle mean [95\% CI] & Learned mean [95\% CI] & Adoption \\
\midrule
Adverse margin & -0.0270 & 0.0000 [0.0000, 0.0000] & -0.0000 [-0.0000, 0.0000] & 0.05\% \\
Null margin & 0.0000 & 0.0000 [0.0000, 0.0000] & -0.0002 [-0.0003, 0.0000] & 5.76\% \\
Low positive dose & 0.0025 & 0.0026 [0.0019, 0.0032] & 0.0001 [-0.0001, 0.0002] & 7.35\% \\
Powered positive dose & 0.0540 & 0.0541 [0.0529, 0.0554] & 0.0512 [0.0500, 0.0525] & 94.17\% \\
\bottomrule
\end{tabular}
}
\caption{Calibrated closure-option mechanism experiment, with 4096 synthetic worlds per prespecified regime. Current frontier and burden are equal by construction; only the innovation-safe library can admit the designated bridge descendant. Intervals are registered finite-simulation uncertainty summaries. This is synthetic mechanism evidence, not a market-return, alpha, forecasting, or causal estimate.}
\label{tab:closure-option-mechanism}
\end{table}


The low-dose oracle has positive value, but the learned-value interval crosses
zero and adoption remains low. At the powered dose, both oracle and learned
mean intervals are positive and learned adoption exceeds the registered gate.
The null point adoption rate passes its prespecified finite-test tolerance, but
its Wilson upper endpoint slightly exceeds that tolerance; both facts are
reported. Adverse adoption is near zero, and every world passes the structural
equal-frontier, equal-burden, and bridge-availability assertions.

\subsubsection{Evidence status and interpretation}

The mechanism result is exact with respect to its structural fixture and
floating-point synthetic evidence with respect to its simulated value and
uncertainty summaries. The complete world ledger is independently regenerated,
its record hashes are checked, and the summary is reproduced. These checks are
not formal proof of a market effect. The powered margin is a prospective power
dose, not an estimated return premium, and the calibrated noise distribution
does not transform the simulation into securities-panel evidence.

Taken together, the studies support a bounded managerial conclusion. Safe
compression can preserve a materially larger research-option menu at much
lower declared maintenance burden, but the decision to exercise any option
requires a separate evidence threshold. The retrospective panel shows
disciplined abstention and one adverse exercise; the synthetic experiment shows
that this outcome can coexist with a valid bridge mechanism when small effects
are difficult to distinguish from proposal noise.

\subsection{Licensed-data and provenance boundary}

All public financial tables in this section are generated from sealed
aggregate manifests. Raw CRSP/WRDS rows, security identifiers, and return paths
are excluded. The exact internal experiment identifiers and hash-bound paths
remain in the repository provenance map so that design and result seals can be
verified; publication-facing scientific names are used here. Reproduction of
the source-data stages requires independent licensed access. Regeneration and
audit of every published aggregate table requires only the committed
public-safe artifacts.
